% TOP_PROBLEM: {"id": 3126, "title": "Lovász Path Removal Conjecture", "category_id": 3, "category": {"id": 3, "name": "graph_theory", "display_name": "Graph Theory", "description": "Problems involving graphs, networks, and their properties.", "slug": "graph-theory", "order_index": 3, "created_at": "2026-07-31T15:26:25.671Z"}, "status": "open", "problem_number": "OPG-46629", "set_id": 12}
% FIRST_POSTED: 2026-10-01
% ATTEMPT_STATUS: unresolved
% UnsolvedMath ID 3126; source catalogue code OPG-46629
% !TeX program = lualatex
% GPT-6 Astra Ultra research attempt, 2026-10-01; self-contained partial work.
% Completion estimate: no numerical percentage assigned; full problem unresolved.
\documentclass[11pt,a4paper]{article}
\usepackage[margin=25mm]{geometry}
\usepackage{fontspec}
\setmainfont{lmroman10-regular.otf}[BoldFont=lmroman10-bold.otf,
 ItalicFont=lmroman10-italic.otf,BoldItalicFont=lmroman10-bolditalic.otf]
\usepackage{amsmath,amssymb,mathtools}
\usepackage{unicode-math}
\setmathfont{latinmodern-math.otf}
\AtBeginDocument{\let\setminus\smallsetminus}
\usepackage{xurl,hyperref}
\hypersetup{colorlinks=true,urlcolor=blue}
\setcounter{secnumdepth}{0}
\setlength{\parindent}{0pt}
\setlength{\parskip}{5pt}
\setlength{\emergencystretch}{3em}
\begin{document}
\section*{Lovász Path Removal Conjecture}\label{3126}
{\small UnsolvedMath ID 3126 · OPG-46629 · Graph Theory · OpenGarden}
\subsection{Definitions and mathematical statement}
A finite simple undirected graph is $k$-vertex-connected, for an integer
$k\ge1$, if it has at least $k+1$ vertices and remains connected after
deletion of any set of fewer than $k$ vertices. A path
$P=(v_0,\ldots,v_t)$ is induced if its vertices are distinct and the
only edges of the graph between these vertices are the consecutive
path edges $v_{i-1}v_i$.

The conjecture asks whether there exists a function
$f:\mathbb N\to\mathbb N$ such that, for every $k\ge1$, every
$f(k)$-vertex-connected graph $G$, and every two distinct vertices
$x,y\in V(G)$, there is an induced path $P$ from $x$ to $y$ for which
\[
 G-V(P)\text{ is }k\text{-vertex-connected}.
\]
Here $G-V(P)$ is the induced subgraph on all vertices outside the path.
All vertices of $P$, including both prescribed endpoints, are deleted.
The retained graph must still have at least $k+1$ vertices under the
connectivity convention above.

The function may grow with $k$, but cannot depend on $G$, its order,
or the choice of endpoints. The conjecture asks only that some finite
threshold exists for each $k$; determining the best possible threshold
is a sharper question.

Deleting just the edges of $P$ is a different and weaker requirement,
since that operation retains the path vertices and their other incident
edges. The path must be induced in the original graph, and the conclusion
must work for every prescribed pair, not only for some convenient pair.
\subsection{Short English statement}
Does sufficiently high vertex-connectivity always allow an induced path between any prescribed pair whose vertex deletion leaves a $k$-connected graph?
\subsection{Research attempt}
\paragraph{Scope and process.}
GPT-6 Astra Ultra, 1 October 2026. This attempt keeps the catalogue statement above, checks primary literature, and develops a restricted argument with an adversarial gap audit. The proofs below are the mathematical output of this attempt, not a claim of novelty or external certification.

\paragraph{Literature and attack plan.}
The checked Garden source [S2] records the known small-$k$ cases and carefully distinguishes vertex deletion from the weaker edge-deletion theorem. Kawarabayashi, Lee, Reed and Wollan's primary publication record [R1] confirms an edge-deletion result; it is not used here to assert the stronger target. We pursue short induced paths, proving an exact deletion-budget lemma, a bounded-diameter case, and complete bipartite examples. This identifies why connectivity alone does not finish this particular strategy.

\paragraph{Connectivity after a bounded deletion.}
If $G$ is $r$-vertex-connected and $X\subseteq V(G)$ has $t<r$ vertices, then $G-X$ is $(r-t)$-vertex-connected. First, it has at least $r+1-t$ vertices. Second, removing any set $Y$ of fewer than $r-t$ additional vertices removes fewer than $r$ vertices in total from $G$, leaving a connected graph. Both the order and deletion requirements in the definition are satisfied.

Consequently, if $G$ is $(k+t)$-connected and has an induced $x$--$y$ path on at most $t$ vertices, its deletion leaves a $k$-connected graph. Notice that $t$ counts vertices, including both endpoints; confusing it with the number of edges would lose one unit in the bound.

\paragraph{A bounded-diameter theorem.}
Suppose the diameter of $G$ is at most $h$ and $G$ is $(k+h+1)$-connected. For any distinct $x,y$, a shortest $x$--$y$ path has at most $h$ edges, hence at most $h+1$ vertices. It is induced: an edge joining nonconsecutive path vertices would shortcut it to a shorter path. The deletion lemma proves
\[
 G-V(P)\text{ is }k\text{-connected}. \tag{1}
\]
In particular, for adjacent endpoints the single edge is induced, and $(k+2)$-connectivity suffices for that pair.

There is a concrete dense subclass with $h=2$. If $G$ has order $n$ and minimum degree at least $\lceil(n-1)/2\rceil$, two nonadjacent vertices have neighbour sets in the remaining $n-2$ vertices whose sizes sum to at least $n-1$. They therefore share a neighbour. Hence diameter is at most two, and every such graph which is $(k+3)$-connected satisfies the target for all prescribed pairs.

\paragraph{Exact calculations for complete bipartite graphs.}
Let $G=K_{a,b}$ with $a,b\ge k+2$. Its vertex-connectivity is $\min(a,b)$: deleting fewer vertices leaves at least one vertex in each part, hence a connected graph; deleting a whole smaller part disconnects the other part. If $x,y$ lie in different parts, take the edge $xy$. The remainder is $K_{a-1,b-1}$ and is at least $(k+1)$-connected.

If $x,y$ lie in the part of size $a$, choose any $z$ in the other part and use $x,z,y$. This path is induced because $x,y$ are nonadjacent. Its deletion leaves $K_{a-2,b-1}$, which is $k$-connected; the case with endpoints in the other part is symmetric. Thus the theorem holds on this family already when both parts have size at least $k+2$, stronger than what the generic diameter budget alone gives.

\paragraph{A necessary lower threshold.}
The complete graph $K_{k+2}$ is $(k+1)$-connected. Every induced path between distinct prescribed endpoints is their single edge, since any longer path has a chord. Removing its two vertices leaves $K_k$, which is not $k$-connected under the required order convention. Thus any universal threshold must satisfy $f(k)\ge k+2$. Complete graphs $K_n$ with $n\ge k+3$ meet the conclusion, so this obstruction is exact within the complete-graph family.

\paragraph{Verification and the gap.}
All constructed paths are induced in the original graph, and the order condition for the residual graph was checked as part of the deletion lemma. In the bipartite calculation no part becomes empty, since $a,b\ge k+2$ and $k\ge1$. The dense degree condition is supplementary; high vertex-connectivity alone was never asserted to imply diameter two.

For a general highly connected graph, shortest paths can still have many vertices relative to a fixed connectivity budget. The inequality ``remaining connectivity is at least initial connectivity minus the path order'' then supplies no useful $k$-bound. Long paths need not destroy connectivity, so failure of this estimate is not a counterexample. The missing structural argument would choose a path whose many vertices can be removed without causing the worst-case loss at each deletion.

\paragraph{Final status.}
The deletion-budget theorem, bounded-diameter subclass, dense subclass, bipartite constructions, and lower threshold are proved. No function depending on $k$ alone has been established for arbitrary graphs. The full conjecture is \textbf{UNRESOLVED} by this attempt.

\subsection{Sources}
\begin{itemize}
\item[S1] ULAM.AI and the UnsolvedMath Contributors, supplied problems.json, record 3126 (OPG-46629), OpenGarden collection. Catalogue identity and source transcription; CC BY 4.0.
\url{https://huggingface.co/datasets/ulamai/UnsolvedMath}.
\item[S2] Open Problem Garden, Lovász Path Removal Conjecture. Original problem and discussion; the mathematical formulation below is expanded from the supplied source record.
\url{http://www.openproblemgarden.org/op/lovasz_path_removal_conjecture}.
\item[R1] K. Kawarabayashi, O. Lee, B. Reed and P. Wollan, \emph{A weaker version of Lovasz' path removal conjecture}, Journal of Combinatorial Theory, Series B 98 (2008), 972--979, DOI 10.1016/j.jctb.2007.11.003. Primary author-repository abstract checked; it concerns edge removal.
\url{https://iris.uniroma1.it/handle/11573/443291}.
\end{itemize}
\end{document}
